% Part: first-order-logic
% Chapter: models-theories
% Section: expressing-props-of-structures

\documentclass[../../../include/open-logic-section]{subfiles}

\begin{document}

\olfileid[pt]{fol}{mat}{exs}

\olsection{Expressando Propriedades de \printtoken{P}{structure}}

\begin{explain}
Frequentemente é útil e importante expressar condições sobre
funções e relações, ou, mais geralmente, que as funções e
relações em uma estrutura satisfaçam essas condições.  Por exemplo, nós
gostaríamos de ter maneiras de distinguir aquelas !!{structure}s para uma
linguagem que ``capturam'' o que queremos que os !!{predicate}s ``signifiquem''
daquelas que não capturam.  É claro que somos completamente livres para especificar
quais !!{structure}s ``pretendemos'', por exemplo, podemos especificar que a
interpretação do !!{predicate}~$\le$ deve ser uma ordenação, ou que
estamos interessados apenas em interpretações de~$\Lang L$ nas quais o
domínio consiste em conjuntos e $\Obj \in$ é interpretado pela relação ``é
!!a{element} de''.  Mas podemos fazer isso com !!{sentence}s da
linguagem?  Em outras palavras, quais condições sobre
!!a{structure}~$\Struct M$ podemos expressar por !!a{sentence} (ou talvez
um conjunto de !!{sentence}s) na linguagem de~$\Struct M$?  Há
algumas condições que não conseguiremos expressar.  Por exemplo,
não há sentença de~$\Lang L_A$ que seja verdadeira em uma
!!{structure}~$\Struct M$ somente se $\Domain M = \Nat$.  Não podemos expressar
``o domínio contém apenas números naturais''.  Mas há
``propriedades estruturais'' de !!{structure}s que talvez possamos
expressar.  Quais propriedades de !!{structure}s podemos expressar por
!!{sentence}s?  Ou, dito de outra forma, quais coleções de
!!{structure}s podemos descrever como aquelas que tornam !!a{sentence} (ou conjunto
de !!{sentence}s) verdadeira?
\end{explain}

\begin{defn}[Modelo de um conjunto]
Seja $\Gamma$ um conjunto de !!{sentence}s em uma linguagem~$\Lang L$.  Nós
dizemos que !!a{structure}~$\Struct M$ \emph{é um modelo de}~$\Gamma$ se
$\Sat{M}{!A}$ para todo $!A \in \Gamma$.
\end{defn}

\begin{ex}
A sentença $\lforall[x][x \le x]$ é verdadeira em~$\Struct M$ sse
$\Assign{\le}{M}$ é uma relação reflexiva.  A sentença
$\lforall[x][\lforall[y][((x \le y \land y \le x) \lif x = y)]]$ é
verdadeira em~$\Struct M$ sse $\Assign{\le}{M}$ é antissimétrica.  A
sentença $\lforall[x][\lforall[y][\lforall[z][((x \le y \land y \le z)
      \lif x \le z)]]]$ é verdadeira em~$\Struct M$ sse $\Assign{\le}{M}$
é transitiva.  Assim, os modelos de
\begin{align*}
\{\quad &\lforall[x][x \le x], \\
   & \lforall[x][\lforall[y][((x \le y \land y \le
    x) \lif x = y)]], \\
   &\lforall[x][\lforall[y][\lforall[z][((x \le y
      \land y \le z) \lif x \le z)]]] \quad \}
\end{align*}
são exatamente aquelas estruturas nas quais~$\Assign{\le}{M}$ é reflexiva,
antissimétrica e transitiva, isto é, uma ordem parcial.  Portanto, podemos tomá-las
como axiomas para a \emph{teoria de primeira ordem das ordens parciais}.
\end{ex}

\end{document}
